\documentclass[10pt,a4paper]{amsart}
\usepackage[margin=19mm]{geometry}
\usepackage{amsmath,amssymb,mathtools}
\usepackage[scr=rsfs,cal=euler]{mathalfa}
\usepackage{xfrac}
\usepackage[unicode,hidelinks,bookmarks=false]{hyperref}
\newcommand{\ie}{i.e.\ }
\newcommand{\cf}{cf.\ }
\newcommand{\resp}{respectively\ }
\newcommand{\Subsection}{\subsection}
\providecommand{\nobreakdash}{\nobreak\hbox{-}}
\setlength{\emergencystretch}{2em}
\multlinegap0pt
\usepackage{bm}
\usepackage{bbm}
\usepackage{dsfont}
\usepackage{xspace}
\usepackage{cancel}
\usepackage{mathtools}
\usepackage{amsthm}
\makeatletter
\@ifundefined{corollary}{\newtheorem{corollary}{Corollary}}{}
\@ifundefined{theorem}{\newtheorem{theorem}{Theorem}}{}
\makeatother
\title[Advanced Distribution Theory for Significance in Scale Space]{Advanced Distribution Theory for Significance in Scale Space}
\author{Rui Liu, Jan Hannig, J. S. Marron}
\date{Source: arXiv:2510.27023v1 (CC BY 4.0)}
\begin{document}
\maketitle

\noindent\emph{Source and licence.} Advanced Distribution Theory for Significance in Scale Space. Version arXiv:2510.27023v1 (CC BY 4.0), distributed under the Creative Commons Attribution 4.0 International licence, as recorded in the arXiv licence metadata for this exact version and shown on the abstract page https://arxiv.org/abs/2510.27023v1. Full rights notice: licenses/arxiv-2510.27023-CC-BY-4.0.txt.\par\smallskip

\noindent\textit{Editorial scope.} Counted unit: the Corollary whose source label is \texttt{corollary-1st-max} in the article (source0.tex lines 287--304, source label \texttt{corollary-1st-max}), delivered together with the article's own complete original proof of it (source0.tex lines 305--347, one \texttt{proof} environment of 43 lines). No mathematics is written, repaired or re-derived: statement and proof are the article's own text line for line. Required context retained uncounted: SSS\_extreme\_1st (lines 255--272); eq\_corr\_g\_1st (lines 258--261). Every definition, display and label of those lines is retained unchanged. Omitted from this package and not redistributed: 1--254, 262--286, 348--1192. Nothing inside a retained excerpt is omitted or altered: every retained body is delivered exactly as the article's lines are written, so no omission receipt of any kind accompanies this package. Citations: the retained excerpts carry no citation command at all, so no bibliography entry of the article is transcribed and no reference is redistributed. Reference adaptation: none was needed and none was made; every reference inside the delivered unit resolves inside it, so no unresolved reference or citation is produced. Printed slips kept literal: no printed word, symbol, hypothesis or number of the counted statement or of its proof was altered. Layout and class: the article's own class is \texttt{article} with packages including graphicx, hyperref, amsgen, amsmath, amstext, amsbsy, amsopn, amssymb, multirow, xspace, epstopdf, authblk, xcolor, xr-hyper; the wrapper is the lane's pinned \texttt{amsart} editorial wrapper at 10pt on A4, which restates the theorem-like environments the retained text needs and adds the article's own macro definitions that the retained text needs (none), with extra wrapper packages (bm, bbm, dsfont, xspace, cancel, mathtools). The delivered numbering is the unit's own. Assets: no retained span contains a figure environment, a graphics-inclusion command, a PDF-page inclusion, a table or a TikZ picture; no image file is copied into this package, the package declares no asset dependency and no image file is redistributed with it. Licence: version arXiv:2510.27023v1 of this article is distributed under the Creative Commons Attribution 4.0 International licence, as recorded in the arXiv licence metadata for this exact version and shown on the abstract page https://arxiv.org/abs/2510.27023v1. A full rights notice is delivered at licenses/arxiv-2510.27023-CC-BY-4.0.txt. Characters: every retained excerpt is pure ASCII, so the delivered engine, XeLaTeX with its default Latin Modern text font, reports no missing character.
\begin{theorem}
	\label{SSS_extreme_1st}
	Consider a  mean-0, variance-1 Gaussian random fields $\hat{T}_{g,i,j}$ indexed by $i,j$. If for each fixed $g$ the random field $\{\hat{T}_{g,i,j}:1\leq i,j\leq g\}$  is stationary with correlation
	\begin{equation}
		\label{eq_corr_g_1st}
		\rho_{g,i,j}=[1-\frac{C^2}{2\log g}(iu+jv)^2]\exp(-\frac{(i^2+j^2)C^2}{4\log g}),
	\end{equation}
	where $C>0$ and $u^2+v^2=1$. Then
	\begin{equation*}
		\lim_{g\rightarrow\infty}P\left\{\max_{i,j=1,...,g}\hat{T}_{g,i,j}\leq u_{g^2}(x)\right\}=\exp(-\vartheta \exp(-x)),
	\end{equation*}
	where 
	\begin{equation}
		\label{eq-vartheta-1st}
		\vartheta\leq2\Phi(C)-1.
	\end{equation}
	%and $u(x)=u_{g^2}(x)$ which is defined in Section \ref{2d_extrem}.
\end{theorem}

\begin{corollary}
	\label{corollary-1st-max}
	Consider two mean-zero, variance-one Gaussian random fields  
	\(\hat{T}^{(1)}_{g,i,j}\) and \(\hat{T}^{(2)}_{g,i,j}\), each indexed by \(1 \le i,j \le g\) and satisfying
	assumption~\eqref{eq_corr_g_1st}. Let \(u_{g^2}(x)\) be the threshold corresponding to marginal tail probability \(\alpha/2\). Then for every fixed \(x \in \mathbb{R}\),
	\[
	\liminf_{g \to \infty}
	P\left\{
	\max_{i,j = 1,\dots, g}
	\left[
	(\hat{T}^{(1)}_{g,i,j})^2 \vee (\hat{T}^{(2)}_{g,i,j})^2
	\right]
	\le u_{g^2}(x)^2
	\right\}
	\ge 4\exp(-\vartheta \exp(-x))-3,
	\]
	where the right-hand side reflects Bonferroni control of the family-wise error rate over the two directional tests.
\end{corollary}

\begin{proof}
	Let us define
	\[
	M^{(1)}_g = \max_{i,j} |\hat{T}^{(1)}_{g,i,j}|, \quad
	M^{(2)}_g = \max_{i,j} |\hat{T}^{(2)}_{g,i,j}|.
	\]
	We are interested in bounding the probability
	\[
	P\left\{
	\max_{i,j} \left[(\hat{T}^{(1)}_{g,i,j})^2 \vee (\hat{T}^{(2)}_{g,i,j})^2\right]
	\le u_{g^2}(x)^2
	\right\}
	=
	P\left\{
	M^{(1)}_g \le u_{g^2}(x), \; M^{(2)}_g \le u_{g^2}(x)
	\right\}.
	\]
	
	Applying the union bound (Bonferroni inequality), we have:
	\[
	P\left\{
	M^{(1)}_g \le u_{g^2}(x), \; M^{(2)}_g \le u_{g^2}(x)
	\right\}
	\ge
	1 - P\left\{M^{(1)}_g > u_{g^2}(x) \right\}
	- P\left\{M^{(2)}_g > u_{g^2}(x) \right\}.
	\]
	
	Under the assumptions of Theorem~\ref{SSS_extreme_1st}, the distribution of \(M^{(k)}_g\) satisfies
	\[
	\lim_{g \to \infty} P\left\{M^{(k)}_g > u_{g^2}(x) \right\} = 2-2\exp(-\vartheta \exp(-x)) \quad (k=1,2).
	\]
	Substituting back into the union bound:
	\[
	\liminf_{g \to \infty} P\left\{
	\max_{i,j} \left[(\hat{T}^{(1)}_{g,i,j})^2 \vee (\hat{T}^{(2)}_{g,i,j})^2\right]
	\le u_{g^2}(x)^2
	\right\}
	\ge 1 - 2 \cdot (2-2\exp(-\vartheta \exp(-x))) = 4\exp(-\vartheta \exp(-x))-3.
	\]
	
	This completes the proof.
\end{proof}

\end{document}
